% Generated draft: requires mathematical and source review.
\documentclass[11pt]{article}
\usepackage[margin=25mm]{geometry}
\usepackage{fontspec}
\setmainfont{Latin Modern Roman}
\usepackage{amsmath,amssymb,mathtools,unicode-math}
\setmathfont{Latin Modern Math}
\usepackage{enumitem,needspace,xurl,hyperref}
\setlength{\parindent}{0pt}
\setlength{\parskip}{5pt}
\setlength{\emergencystretch}{3em}
\newcommand{\N}{\mathbb{N}}
\newcommand{\Z}{\mathbb{Z}}
\newcommand{\Q}{\mathbb{Q}}
\newcommand{\R}{\mathbb{R}}
\newcommand{\C}{\mathbb{C}}
\newcommand{\F}{\mathbb{F}}
\newcommand{\A}{\mathbb{A}}
\newcommand{\PP}{\mathbb{P}}
\newcommand{\E}{\mathbb{E}}
\newcommand{\eps}{\varepsilon}
\newcommand{\dd}{\,\mathrm{d}}
\begin{document}
\section*{UnsolvedMath ID: 2308022}
\textbf{Source problem code:} catalogue entry
\subsection{Definitions and mathematical statement}
For a bounded smoothly bounded domain $\Omega\subset\C^n$, a standard inner-function convention is a bounded holomorphic function $f:\Omega\to\mathbb D$ whose non-tangential boundary values satisfy $|f^*(\zeta)|=1$ for surface-almost every $\zeta\in\partial\Omega$. The intended function must be nonconstant; otherwise unimodular constants give a trivial answer under definitions that allow them. This convention requires a suitable boundary and boundary-value measure, not specified in the source.

Under a $C^2$ defining-function convention, strict pseudoconvexity means that locally $\Omega=\{\rho<0\}$, $d\rho\ne0$ on the boundary, and
\[
 \sum_{j,k=1}^n\frac{\partial^2\rho}{\partial z_j\partial\bar z_k}(\zeta)
               v_j\overline{v_k}>0
\]
for every nonzero complex tangent vector $v$ at each boundary point $\zeta$. The source does not state boundedness or boundary regularity, so this supplies a standard version rather than silently narrowing ``any strictly pseudoconvex domain.''

\textbf{Question (R. Zeinstra).} For every strictly pseudoconvex domain $\Omega\subset\C^n$, $n\geq2$, in the intended domain and boundary-value class, does there exist a nonconstant inner function? In the smooth bounded version this asks for
\[
 \exists f\in\mathcal O(\Omega),\quad |f|<1\text{ on }\Omega,\quad
 f\text{ nonconstant},\quad |f^*|=1\text{ surface-almost everywhere on }\partial\Omega.
\]
The universal quantifier is over all domains of the intended type, not merely balls or one favorable subclass.

The source reports existence for the ball, attributed to Alexandrov and independently to L\o w and to Hakim--Sibony. The research summary likewise marks the ball solved (mentioning $n\geq1$ for that result), reports positive results for further domains under additional conditions, and leaves the fully general strictly pseudoconvex case open. Its mixed ``SOLVED-IN-LITERATURE'' and partial-progress wording is interpreted only at those stated scopes; it supplies neither a precise general-domain theorem nor hypotheses for the further positive cases. No scope is enlarged here.
\subsection{Short English statement}
Does every strictly pseudoconvex domain in $\C^n$, $n\geq2$, admit a nonconstant holomorphic inner function? The record reports affirmative results for the ball and some other classes while leaving the fully general question open. Its boundary regularity and boundary-value conventions are unspecified.
\subsection{Sources}
\begin{itemize}
\item Local UnsolvedMath registry, permanent ID 2308022, source code catalogue entry, ``Research Problems in Function Theory --- Problem 8.22''; attribution: R. Zeinstra. Supplied status partially solved, classification PARTIAL-PROGRESS, updated 2026-08-06; the summary distinguishes the solved ball case from the general question.
\item Hayman, \emph{Research Problems in Function Theory} (2018), Problem 8.22, \url{https://arxiv.org/abs/1809.07200}. The record reports source-TeX extraction and access on 2026-07-29; the 2018 update was retained for status triage. Current status and rights require review; the record describes publicly available arXiv source TeX.
\item The source names Alexandrov, L\o w, and Hakim--Sibony for ball existence; it expressly gives no citation for Alexandrov or L\o w and supplies no complete Hakim--Sibony reference. The summary mentions extensions of Alexandrov's method under additional conditions without identifying papers or exact hypotheses.
\end{itemize}
\end{document}
